
\subsubsection{3.1 Formality Measurement}

Formality is measured using the DeBERTa Large Formality Ranker (s-nlp/deberta-large-formality-ranker), a transformer model fine-tuned on the GYAFC dataset with 87.8\% classification accuracy. The model produces continuous scores in [-1, 1], where higher values indicate greater formality.

For each conversation turn, the text is extracted and its formality score computed. The formality delta between human input and AI response is:

$$\Delta = |f_{\text{AI}_1} - f_{\text{H}_1}|$$

Lower delta indicates stronger accommodation---the AI matches the human's formality level more closely.

\subsubsection{3.2 Dataset and Preprocessing}

The study uses the Anthropic hh-rlhf dataset, which contains over 170,000 human-AI conversations with clear turn structure. The following filters are applied:

\begin{itemize}
\item \textbf{Minimum turns}: $\geq 2$ turns per participant (for H-E1, H-M1: $\geq 4$ turns per role for trajectory analysis)
\item \textbf{Language}: English only
\item \textbf{Message length}: $\geq 5$ tokens
\end{itemize}

After filtering, 111,039 conversations are retained for primary correlation analysis (H-M2), and 26,395 conversations for trajectory-based analyses requiring longer exchanges (H-E1, H-M1).

\subsubsection{3.3 Analysis Framework}

\textbf{H-E1 (Existence)}: Bidirectional Convergence Scores (BCS) are computed as Pearson correlation between user and AI complexity trajectories within each conversation. Success criterion: SD > 0.15, n > 10,000.

\textbf{H-M1 (User Adaptation)}: Lagged cross-correlation is computed between AI complexity at turn t and user complexity at turn t+1. Success criterion: r > 0, p < 0.05.

\textbf{H-M2 (AI Accommodation)}: Pearson correlation is computed between human formality and AI formality in turn-1 pairs. Success criterion: |r| > 0.1, p < 0.001.

\textbf{H-M3 (Engagement Link)}: Conversations are divided into terciles by formality delta magnitude. Continuation rates are measured per tercile. Original prediction: monotonic trend T1 > T2 > T3 (lowest delta has highest continuation).

